% =====================================================================
\chapter{Derivadas parciales y gradiente}
\label{tema:1}
% =====================================================================

La pérdida de un modelo de lenguaje es una función escalar de muchas variables. Este tema define las derivadas parciales y la diferenciabilidad, introduce el gradiente y la derivada direccional, y demuestra que el gradiente señala la dirección de máximo ascenso, que es el fundamento del descenso de gradiente. Se siguen las clases~7, 10 y~11 del curso 18.02 del MIT~\cite{mit1802}, la §5.2 de~\cite{mml} y la §2.4 de~\cite{d2l}.

\section{Derivadas parciales}

Sea $U\subseteq\R^{n}$ un abierto y $f\colon U\to\R$. Para estudiar cómo varía $f$ cerca de un punto, se empieza por moverse en una sola coordenada.

\begin{definicion}[Derivada parcial]\label{def:parcial}
La \emph{derivada parcial} de $f$ respecto de $x_i$ en el punto $x\in U$ es
\[
\frac{\partial f}{\partial x_i}(x)=\lim_{h\to 0}\frac{f(x+h\,e_i)-f(x)}{h},
\]
si el límite existe. Equivale a derivar $f$ como función de una sola variable $x_i$, tratando las demás como constantes.
\end{definicion}

\begin{ejemplo}\label{ej:parciales}
Sea $f(x,y)=x^{2}y+3y$. Con $y$ fija, $\partial f/\partial x=2xy$; con $x$ fija, $\partial f/\partial y=x^{2}+3$. En el punto $(2,1)$, donde $f(2,1)=7$, se obtiene $\partial f/\partial x=4$ y $\partial f/\partial y=7$: al aumentar ligeramente $x$, la función crece a razón de~4 unidades por unidad de $x$; al aumentar $y$, a razón de~7.
\end{ejemplo}

\begin{observacion}[Conexión con los LLM]
Si $L(\theta)$ es la pérdida de entropía cruzada de un modelo con parámetros $\theta_1,\dots,\theta_N$, la parcial $\partial L/\partial\theta_i$ mide cuánto aumenta la pérdida, a primer orden, al aumentar el parámetro $\theta_i$ con los demás fijos. Cada paso de entrenamiento calcula las $N$ parciales a la vez.
\end{observacion}

\section{Diferenciabilidad y gradiente}

La existencia de las derivadas parciales no basta para que $f$ se comporte bien cerca de $x$: las parciales solo exploran las direcciones de los ejes. La noción adecuada exige que $f$ admita una buena aproximación lineal en todas las direcciones a la vez~\cite[§5.2]{mml}.

\begin{definicion}[Diferenciabilidad y gradiente]\label{def:diferenciable}
La función $f$ es \emph{diferenciable} en $x\in U$ si existe un vector $g\in\R^{n}$ tal que
\[
f(x+h)=f(x)+g\T h+o(\norm{h})\qquad\text{cuando } h\to 0,
\]
es decir, $\lim_{h\to 0}\bigl(f(x+h)-f(x)-g\T h\bigr)/\norm{h}=0$. El vector $g$ se llama \emph{gradiente} de $f$ en $x$ y se escribe $\nabla f(x)$. La función afín $h\mapsto f(x)+\nabla f(x)\T h$ es la \emph{aproximación lineal} (o de primer orden) de $f$ en $x$.
\end{definicion}

\begin{proposicion}\label{prop:grad-parciales}
Si $f$ es diferenciable en $x$, existen todas las parciales en $x$ y
\[
\nabla f(x)=\Bigl(\frac{\partial f}{\partial x_1}(x),\dots,\frac{\partial f}{\partial x_n}(x)\Bigr)\T.
\]
En particular, el gradiente es único.
\end{proposicion}

\begin{proof}
Sea $g$ un vector como en la \cref{def:diferenciable}. Tomando $h=t\,e_i$ con $t\to 0$, se tiene $\norm{h}=\abs{t}$ y $g\T h=t\,g_i$, luego
\[
\frac{f(x+t\,e_i)-f(x)}{t}-g_i=\frac{f(x+t\,e_i)-f(x)-g\T h}{t}\longrightarrow 0,
\]
porque el cociente por $t$ difiere del cociente por $\abs{t}$ solo en el signo. Así, $\partial f/\partial x_i(x)$ existe y vale $g_i$. Como las parciales están determinadas por $f$, también lo está $g$.
\end{proof}

El gradiente tiene la misma forma que la variable: si $x\in\R^{n}$, entonces $\nabla f(x)\in\R^{n}$. En el \cref{ej:parciales}, $\nabla f(2,1)=(4,7)\T$. Para un LLM, el gradiente de la pérdida tiene tantas componentes como parámetros.

\begin{teorema}[Condición suficiente de diferenciabilidad]\label{teo:C1}
Si todas las derivadas parciales de $f$ existen en un entorno de $x$ y son continuas en $x$, entonces $f$ es diferenciable en $x$. En particular, toda $f\in C^{1}(U)$ es diferenciable en todo punto de $U$.
\end{teorema}

La demostración, estándar en los cursos de cálculo de varias variables, aplica el teorema del valor medio en cada coordenada por separado y se omite aquí. Las funciones que se usan para construir una red (polinomios, exponenciales, logaritmos, cocientes con denominador no nulo y sus composiciones) son de clase $C^{\infty}$ en su dominio. La excepción habitual es la función ReLU, $\max(0,t)$, que no es derivable en $t=0$; el \cref{tema:6} vuelve sobre ella.

\begin{observacion}[Parciales sin diferenciabilidad]
Sea $f(x,y)=xy/(x^{2}+y^{2})$ si $(x,y)\ne(0,0)$ y $f(0,0)=0$. Como $f$ se anula sobre los dos ejes, $\partial f/\partial x(0,0)=\partial f/\partial y(0,0)=0$. Sin embargo, $f(t,t)=1/2$ para todo $t\ne 0$, de modo que $f$ ni siquiera es continua en el origen y, por tanto, no es diferenciable allí. Las parciales existen, pero no son continuas en $(0,0)$.
\end{observacion}

\section{Derivada direccional}

\begin{definicion}[Derivada direccional]\label{def:direccional}
Sea $u\in\R^{n}$ con $\norm{u}=1$. La \emph{derivada direccional} de $f$ en $x$ en la dirección $u$ es
\[
D_u f(x)=\lim_{t\to 0}\frac{f(x+t\,u)-f(x)}{t}.
\]
Las derivadas parciales son las derivadas direccionales en las direcciones $e_1,\dots,e_n$.
\end{definicion}

\begin{proposicion}\label{prop:direccional}
Si $f$ es diferenciable en $x$, entonces $D_u f(x)=\nabla f(x)\T u$ para todo $u$ unitario.
\end{proposicion}

\begin{proof}
Con $h=t\,u$ se tiene $\norm{h}=\abs{t}$. Por la \cref{def:diferenciable}, $f(x+t\,u)-f(x)=t\,\nabla f(x)\T u+o(\abs{t})$. Dividiendo por $t$ y haciendo $t\to 0$ se obtiene el enunciado.
\end{proof}

\begin{ejemplo}\label{ej:direccional}
Con la función del \cref{ej:parciales} y $u=(3/5,\,4/5)\T$, que es unitario, $D_u f(2,1)=4\cdot 0{,}6+7\cdot 0{,}8=8$. Se verá a continuación que ninguna dirección supera $\norm{\nabla f(2,1)}=\sqrt{65}\approx 8{,}062$.
\end{ejemplo}

\section{La dirección de máximo ascenso}

\begin{teorema}[Máximo ascenso]\label{teo:max-ascenso}
Sea $f$ diferenciable en $x$ con $\nabla f(x)\ne 0$. Para todo $u$ unitario,
\[
-\norm{\nabla f(x)}\le D_u f(x)\le\norm{\nabla f(x)}.
\]
La cota superior se alcanza únicamente en $u^{*}=\nabla f(x)/\norm{\nabla f(x)}$ y la inferior únicamente en $-u^{*}$.
\end{teorema}

\begin{proof}
Por la \cref{prop:direccional} y la desigualdad de Cauchy--Schwarz (módulo~2),
\[
\abs{D_u f(x)}=\abs{\nabla f(x)\T u}\le\norm{\nabla f(x)}\,\norm{u}=\norm{\nabla f(x)}.
\]
La igualdad en Cauchy--Schwarz se da si y solo si $u$ es proporcional a $\nabla f(x)$; al ser $u$ unitario, $u=\pm u^{*}$. Sustituyendo, $\nabla f(x)\T u^{*}=\norm{\nabla f(x)}^{2}/\norm{\nabla f(x)}=\norm{\nabla f(x)}$, y con $-u^{*}$ se obtiene el valor opuesto.
\end{proof}

Así, $\nabla f(x)$ apunta en la dirección en la que $f$ crece más deprisa, su norma es esa tasa máxima de crecimiento y $-\nabla f(x)$ es la dirección de \emph{máximo descenso}.

\begin{lema}[Derivada a lo largo de una curva]\label{lem:curva}
Sea $r\colon(a,b)\to U$ derivable en $t_0$ y $f$ diferenciable en $r(t_0)$. Entonces $t\mapsto f(r(t))$ es derivable en $t_0$ y
\[
\frac{d}{dt}f(r(t))\Big|_{t=t_0}=\nabla f(r(t_0))\T r'(t_0).
\]
\end{lema}

\begin{proof}
Sea $x=r(t_0)$ y $h(s)=r(t_0+s)-x$, de modo que $h(s)=s\,r'(t_0)+o(s)$ y $\norm{h(s)}=O(\abs{s})$. Por la \cref{def:diferenciable}, $f(r(t_0+s))-f(x)=\nabla f(x)\T h(s)+o(\norm{h(s)})=s\,\nabla f(x)\T r'(t_0)+o(s)$. Dividiendo por $s$ y haciendo $s\to 0$ se obtiene el enunciado.
\end{proof}

El \cref{lem:curva} es un caso particular de la regla de la cadena, que se estudia en general en el \cref{tema:2}.

\begin{corolario}[Gradiente y conjuntos de nivel]\label{cor:nivel}
Sea $r$ una curva derivable contenida en un conjunto de nivel de $f$, es decir, $f(r(t))=c$ para todo $t$. Entonces $\nabla f(r(t))\T r'(t)=0$: el gradiente es ortogonal a la recta tangente a la curva de nivel.
\end{corolario}

\begin{proof}
La función $t\mapsto f(r(t))$ es constante; su derivada es nula y, por el \cref{lem:curva}, es igual a $\nabla f(r(t))\T r'(t)$.
\end{proof}

\begin{figure}[h]
\centering
\begin{tikzpicture}[scale=1.45,font=\small]
  \draw[->,gray] (-2.4,0) -- (2.8,0) node[right] {$x$};
  \draw[->,gray] (0,-1.4) -- (0,2.3) node[above] {$y$};
  \foreach \c in {1,2.5,4} {
    \draw[thick] (0,0) ellipse ({sqrt(\c)} and {sqrt(\c/3)});
  }
  \draw[dashed,gray!60!black] (-0.65,1.55) -- (2.35,0.55);
  \draw[->,>={Stealth[length=2.2mm]},very thick] (1,1) -- (1.33,1.99);
  \node[right] at (1.36,1.95) {$\nabla f(1,1)=(2,6)\T$};
  \fill (1,1) circle (1.4pt);
  \node[below left] at (1,1) {$(1,1)$};
  \node at (-1.75,1.25) {$f=4$};
\end{tikzpicture}
\caption{Curvas de nivel $x^{2}+3y^{2}=c$ para $c=1$; $2{,}5$; $4$. En $(1,1)$ el gradiente (flecha, a escala reducida) es perpendicular a la tangente a la curva $f=4$ (trazo discontinuo), de dirección $(3,-1)\T$.}
\label{fig:nivel}
\end{figure}

\section{Descenso de gradiente}

El \cref{teo:max-ascenso} sugiere un método para minimizar $f$: desplazarse repetidamente en la dirección $-\nabla f$.

\begin{definicion}[Descenso de gradiente]\label{def:dg}
Dado un punto inicial $x_0$ y una \emph{tasa de aprendizaje} $\eta>0$, el \emph{descenso de gradiente} genera la sucesión
\[
x_{k+1}=x_k-\eta\,\nabla f(x_k),\qquad k=0,1,2,\dots
\]
\end{definicion}

\begin{proposicion}[Descenso a primer orden]\label{prop:descenso}
Si $f$ es diferenciable en $x$ y $g=\nabla f(x)\ne 0$, existe $\eta_0>0$ tal que $f(x-\eta g)<f(x)$ para todo $\eta\in(0,\eta_0)$.
\end{proposicion}

\begin{proof}
Por la \cref{def:diferenciable} con $h=-\eta g$, $f(x-\eta g)=f(x)-\eta\norm{g}^{2}+o(\eta)$. Luego $\bigl(f(x-\eta g)-f(x)\bigr)/\eta\to-\norm{g}^{2}<0$ cuando $\eta\to 0^{+}$, y el cociente es negativo para todo $\eta$ suficientemente pequeño.
\end{proof}

La proposición garantiza el descenso solo para pasos pequeños. El tamaño máximo admisible depende de la curvatura de $f$ y se determina en el \cref{tema:4}.

\begin{ejemplo}\label{ej:dg}
Sea $f(x,y)=x^{2}+3y^{2}$, con $\nabla f=(2x,6y)\T$. Desde $(1,1)$, con $\nabla f(1,1)=(2,6)\T$ y $\eta=0{,}1$, el primer paso da $(1,1)-0{,}1\,(2,6)=(0{,}8;\,0{,}4)$, y la función baja de $f(1,1)=4$ a $f(0{,}8;\,0{,}4)=0{,}64+0{,}48=1{,}12$.
\end{ejemplo}

\begin{observacion}[Conexión con los LLM]
El entrenamiento de un modelo de lenguaje aplica la \cref{def:dg} a la pérdida $L(\theta)$, con el gradiente estimado sobre un lote de datos. Los métodos que se usan en la práctica (descenso estocástico, momento, Adam) modifican esta iteración y se estudian en el módulo~7.
\end{observacion}

\begin{observacion}[Relación con la física]
Una partícula en un potencial $U(x)$ sometida a un rozamiento muy fuerte se mueve con velocidad proporcional a la fuerza $-\nabla U$: $\dot x(t)=-\nabla U(x(t))$. Esta ecuación se llama \emph{flujo de gradiente}, y la \cref{def:dg} es su discretización por el método de Euler con paso $\eta$. Añadir un término de ruido da la dinámica de Langevin, que sirve de modelo para el descenso de gradiente estocástico.
\end{observacion}

\begin{lecturas}
\item MIT 18.02~\cite{mit1802}, clase~7 «Level Curves, Partial Derivatives, Tangent Plane» y clase~11 «Gradient, Directional Derivative».
\item Deisenroth, Faisal y Ong~\cite{mml}, §5.2 «Partial Differentiation and Gradients».
\item Zhang \emph{et al.}~\cite{d2l}, §2.4.3 «Partial Derivatives and Gradients».
\item Goodfellow, Bengio y Courville~\cite{gbc}, §4.3 «Gradient-Based Optimization».
\end{lecturas}

\section*{Ejercicios del tema 1}
\addcontentsline{toc}{section}{Ejercicios del tema 1}

\begin{ejercicio}{1}{eje:1-vf}
Justifique si son verdaderas o falsas: (a) el gradiente es perpendicular a las curvas de nivel; (b) $-\nabla f$ es la dirección de máximo descenso; (c) si existen todas las derivadas parciales de $f$ en $x$, entonces $f$ es diferenciable en $x$.
\end{ejercicio}

\begin{ejercicio}{2}{eje:1-gradiente}
Sea $f(x,y)=x^{3}+2xy^{2}$. Calcule $\nabla f(1,2)$, su norma y la derivada direccional en la dirección $(1,1)\T/\sqrt{2}$. Compruebe que esta última no supera la norma del gradiente.
\end{ejercicio}

\begin{ejercicio}{2}{eje:1-d2l}
Sea $f(x)=3x_1^{2}+5e^{x_2}$. Calcule $\nabla f$ en un punto genérico y en $(1,0)$. Halle la derivada direccional en $(1,0)$ en la dirección del vector $(3,4)\T$ y compárela con $\norm{\nabla f(1,0)}$.
\fuente{adaptado de Zhang \emph{et al.}~\cite{d2l}, §2.4, ejercicio~7 (se añaden el punto y la dirección).}
\end{ejercicio}

\begin{ejercicio}{2}{eje:1-norma}
Calcule el gradiente de $f(x)=\norm{x}$ para $x\in\R^{n}\setminus\{0\}$ y evalúelo en $x=(3,4)\T$. Estudie qué ocurre en $x=0$. Relacione el resultado con el recorte de gradientes por norma, que divide un gradiente $g$ por $\norm{g}$ cuando esta supera un umbral.
\fuente{adaptado de Zhang \emph{et al.}~\cite{d2l}, §2.4, ejercicio~8.}
\end{ejercicio}

\begin{ejercicio}{2}{eje:1-mit}
Sea $f(x,y)=x^{3}-3x^{2}+xy-y^{2}+x+y$. En el punto $(1,1/2)$, determine la dirección unitaria en la que la derivada direccional es máxima y su valor, la dirección en la que es mínima y las direcciones en las que es nula.
\fuente{adaptado de MIT 18.02, primavera de 2006~\cite{mit1802}, \emph{Problem Set}~4, segunda parte, problema~5.}
\end{ejercicio}

\begin{ejercicio}{3}{eje:1-minimo}
Sea $f$ diferenciable en $x$ con $\nabla f(x)\ne 0$. Demuestre, sin usar el \cref{teo:max-ascenso}, que el mínimo de $D_u f(x)$ sobre los vectores unitarios vale $-\norm{\nabla f(x)}$ y se alcanza solo en $u=-\nabla f(x)/\norm{\nabla f(x)}$. Deduzca, para $f(x,y)=x^{2}+y^{2}$, que el gradiente es perpendicular a las circunferencias $x^{2}+y^{2}=c$ parametrizándolas como $r(t)=\sqrt{c}\,(\cos t,\sin t)\T$.
\end{ejercicio}
